% !TeX program = xelatex
\documentclass[11pt]{article}

% ===== Font =====
\usepackage{fontspec}
\usepackage{xeCJK}

% 英文主字型（論文風）
\setmainfont{Times New Roman}

% 中文字型（乾淨）
\setCJKmainfont{Microsoft JhengHei}

% ===== Math =====
\usepackage{amsmath, amssymb, bm}

% ===== Layout =====
\usepackage[a4paper, margin=1in]{geometry}

% ===== Figures =====
\usepackage{graphicx}
\usepackage{subcaption}

% ===== Tables =====
\usepackage{booktabs}

% ===== Links =====
\usepackage[colorlinks=true, linkcolor=blue, citecolor=blue]{hyperref}

% ===== Better spacing =====
\usepackage{setspace}
\onehalfspacing
% Entropy
\newcommand{\entropy}{H}

% Reward
\newcommand{\R}{R_{\text{behavior}}

}

% Text helpers
\newcommand{\abstain}{\text{abstain}}
\newcommand{\correct}{\text{correct}}
\newcommand{\wrong}{\text{wrong}}

\title{Learning When to Answer: Behavior-Oriented Reinforcement Learning for Hallucination Mitigation}
\author{Max \and Eddie}

\begin{document}

\maketitle

\begin{abstract}
Large language models (LLMs) often hallucinate under uncertainty, producing fluent yet unsupported or incorrect responses. 
We argue that hallucination is fundamentally a decision-making problem, involving when to answer, when to abstain, and how confidently to respond, rather than purely a generation error.

To address this, we propose a behavior-oriented reinforcement learning framework that explicitly models these decisions. 
Our method integrates behavior alignment, entropy-based uncertainty modeling, response quality shaping, and length regularization into a unified reward function. 
The model is optimized through a two-stage training process, combining preference learning via Direct Preference Optimization (DPO)~\cite{dpo} and reinforcement learning with Group Relative Policy Optimization (GRPO)~\cite{GRPO}.

Experimental results show that our approach reduces hallucination by over 65\%, consistently outperforming strong baselines such as GPT-4o. 
Moreover, the learned behavior generalizes effectively to unseen domains, improves response quality across multiple evaluation dimensions, and preserves general knowledge capability without catastrophic forgetting~\cite{EWC}.

These results demonstrate that hallucination mitigation can be effectively achieved by learning behavior under uncertainty, enabling LLMs to produce responses that are both reliable and useful.
\end{abstract}
\section{Introduction}

Large language models (LLMs) have demonstrated remarkable capabilities across a wide range of tasks. 
However, they are prone to hallucination, generating fluent but unsupported or incorrect responses~\cite{ji2023survey,Ji_2023,lin2022truthfulqameasuringmodelsmimic}, particularly under conditions of uncertainty or incomplete information. 
Such behavior poses significant challenges for real-world applications where reliability is critical.

Existing approaches to mitigate hallucination primarily focus on improving knowledge access or post-hoc verification. 
For example, retrieval-augmented generation (RAG)~\cite{rag} enhances factual grounding by incorporating external documents, while various detection methods attempt to identify hallucinated outputs after generation~\cite{lin2022truthfulqameasuringmodelsmimic}. 
Despite their effectiveness, these approaches do not explicitly address a fundamental issue: LLMs often fail to make appropriate decisions under uncertainty, such as when to answer, when to abstain, and how confidently to respond.

In this work, we argue that hallucination is not merely a generation error, but a consequence of suboptimal decision-making under uncertainty. 
From this perspective, mitigating hallucination requires learning a policy that balances correctness, confidence, and usefulness, rather than solely improving token-level prediction.

To this end, we propose a behavior-oriented reinforcement learning framework that explicitly models hallucination as a decision problem. 
Our approach optimizes model behavior along three key dimensions:
(i) whether to answer or abstain, based on the availability of reliable information; 
(ii) how to construct responses, ensuring grounded and informative outputs; and 
(iii) how confident to be, avoiding overconfident generation under uncertainty.

We design a unified reward function that integrates multiple components: 
behavior alignment, which encourages correct decision boundaries between answering and abstaining; 
uncertainty modeling via entropy-based overconfidence penalties; 
response quality shaping to promote completeness, actionability, and the use of supporting evidence; 
and length regularization to prevent degenerate short responses or excessively verbose outputs.

To optimize this objective, we adopt a two-stage training framework consisting of preference-based initialization using Direct Preference Optimization (DPO)~\cite{dpo}, followed by reinforcement learning with Group Relative Policy Optimization (GRPO)~\cite{GRPO}, enabling stable and fine-grained policy optimization.

We conduct extensive experiments to evaluate the proposed method. 
The results show that our approach reduces hallucination by over 65\%, consistently outperforming strong baselines such as GPT-4o. 
Furthermore, the learned behavior generalizes effectively to unseen domains, improves response quality across multiple evaluation dimensions, and preserves general knowledge capability without catastrophic forgetting~\cite{EWC}.

These findings demonstrate that hallucination mitigation can be effectively achieved by modeling it as a behavior learning problem under uncertainty, rather than solely improving knowledge access or generation mechanisms.
\section{Method}

\subsection{Problem Formulation}

We consider hallucination in large language models as a decision-making problem under uncertainty, rather than purely a generation problem.

Given an input query $x$ and contextual information $C = \{c_1, c_2, \dots, c_k\}$, a language model generates a response $y$. 
The objective is to produce responses that are both grounded and useful, while avoiding unsupported or misleading outputs.

\subsubsection{Hallucination as Decision-Making}

We define hallucination as a failure to make appropriate decisions under uncertainty. 
Specifically, the model must jointly determine: 
(i) whether to answer or abstain, 
(ii) how to construct the response, and 
(iii) how confident the response should be.

Under this formulation, hallucination occurs when the model produces responses that are either unsupported by available information or expressed with excessive confidence under uncertainty.

\subsubsection{Controlled Contextual Scenarios}

To systematically study model behavior, we consider two controlled scenarios:

\begin{itemize}
    \item \textbf{No reliable information:} The context does not contain sufficient evidence to answer the query. The desired behavior is to abstain.
    \item \textbf{Partially correct information:} The context contains a mixture of relevant and irrelevant information, including at least one correct supporting evidence. The desired behavior is to generate a selective and grounded response.
\end{itemize}

These scenarios enable explicit modeling of the decision boundary between answering and abstaining.

\subsubsection{Groundedness and Uncertainty}

We introduce a groundedness score $s \in [0,1]$, measuring the degree to which a response is supported by the available context, following the groundedness metric proposed in RAG evaluation frameworks \cite{es2024ragas}.

We quantify model uncertainty using the average token-level entropy, which serves as a proxy for predictive uncertainty and has been widely adopted in prior work~\cite{farquhar2024detecting}:

\[
H_{\text{avg}} = -\frac{1}{T} \sum_{t=1}^{T} \sum_{v \in \mathcal{V}} p_t(v)\log p_t(v)
\]

where $T$ is the response length and $p_t(v)$ denotes the probability of token $v$ at step $t$.

We interpret low entropy as high confidence and high entropy as high uncertainty. 
Hallucination often arises when the model exhibits low entropy despite low groundedness.

\subsubsection{Objective}

Our objective is to learn a policy that:
(i) abstains when reliable information is unavailable, 
(ii) produces grounded responses when sufficient evidence exists, 
(iii) avoids overconfident behavior under uncertainty, and 
(iv) maintains response quality and usefulness.

This formulation naturally leads to a reinforcement learning framework that optimizes behavior under uncertainty rather than purely maximizing likelihood.


\subsection{Reward Function}

We design a behavior-oriented reward function that explicitly models hallucination as a decision-making problem under uncertainty. 
The reward integrates four components: behavior alignment, confidence control, response quality, and length regularization.

\[
R = R_{\text{behavior}} + R_{\text{confidence}} + R_{\text{quality}} + R_{\text{length}}
\]

\subsubsection{Behavior Alignment}

We distinguish two controlled scenarios: (i) no reliable information, where the model should abstain, and (ii) partially correct information, where the model should answer.

Let $a \in \{0,1\}$ denote the model action, where $a=1$ indicates abstention and $a=0$ indicates answering.

\[
R_{\text{behavior}} =
\begin{cases}
+R_{\text{abstain}}, & \text{if abstention is correct} \\
-\alpha, & \text{if abstention is incorrect} \\
0, & \text{otherwise}
\end{cases}
\]

\subsubsection{Confidence Control via Entropy}

We introduce an overconfidence penalty based on entropy:

\[
H_{\text{avg}} = -\frac{1}{T} \sum_{t=1}^{T} \sum_{v \in \mathcal{V}} p_t(v)\log p_t(v)
\]

\[
\text{overconf} = \max(0, U_{\min} - H_{\text{avg}})
\]

\[
R_{\text{confidence}} = -\gamma (1 - s)\cdot \text{overconf}
\]

This discourages overly confident predictions when groundedness is low, while preserving confident generation when responses are well-supported.

\subsubsection{Response Quality Shaping}

We introduce a quality score $q \in [0,1]$, computed from multiple dimensions including completeness, actionability, evidence use, and conciseness, following LLM-based evaluation approaches such as G-Eval~\cite{geval}.

\[
R_{\text{quality}} = \beta \cdot q
\]

This component promotes informative and well-structured responses beyond factual correctness.

\subsubsection{Length Regularization}

To prevent degenerate short responses, we impose a length constraint:

\[
L \in [L_{\text{low}}, \infty]
\]

\[
R_{\text{length}} =
-\eta \left(
\max\left(0, \frac{L_{\text{low}} - L}{L_{\text{low}}}\right)
\right)
\]

This penalizes under-informative outputs. In practice, we set $L_{\text{low}} = 50$.

\subsubsection{Summary}

The reward function jointly optimizes decision-making, uncertainty-aware behavior, response quality, and generation stability, enabling balanced behavior under uncertainty.


\subsection{Training Framework}

We adopt a two-stage LoRA training framework consisting of preference-based initialization followed by reinforcement learning refinement.

\subsubsection{Stage 1: Preference Optimization (DPO)}

We first perform Direct Preference Optimization (DPO)~\cite{dpo} to obtain a stable initialization.

Given $(x, C)$, we construct preference pairs $(y^{+}, y^{-})$:

\begin{itemize}
    \item $y^{+}$: preferred response (grounded or correct)
    \item $y^{-}$: less preferred response (hallucinated or incorrect)
\end{itemize}

\[
\mathcal{L}_{\text{DPO}} = - \log \sigma \left( \beta_{\text{dpo}} \left( \log \pi_{\theta}(y^{+}|x,C) - \log \pi_{\theta}(y^{-}|x,C) \right) \right)
\]

\subsubsection{Stage 2: Reinforcement Learning (GRPO)}

We refine the model using Group Relative Policy Optimization (GRPO)~\cite{GRPO}.

For each $(x, C)$, we sample:

\[
\{y_1, y_2, \dots, y_N\} \sim \pi_{\theta}(\cdot | x, C)
\]

Each response is evaluated using the reward $R(y)$.

\[
A_i = R(y_i) - \frac{1}{N} \sum_{j=1}^{N} R(y_j)
\]

The policy is updated to increase the likelihood of higher-reward responses relative to the group.

\subsubsection{Training Stability}

We observe that reinforcement learning may lead to degenerate strategies such as overly short responses. 
To mitigate this, we incorporate quality-based reward shaping and length regularization to maintain a balance between hallucination mitigation and response usefulness.
\section{Experiments}

\subsection{Experimental Setup}

\subsubsection{Models}

We evaluate our method on a 12B-parameter open-weight language model (Gemma3-12B-IT). 
We compare against the base model (without fine-tuning) and GPT-4o as a strong closed-source baseline.

\subsubsection{Evaluation Settings}

We conduct experiments under controlled contextual conditions with varying information quality, including scenarios with no reliable information and scenarios with partially correct information. 
This setup enables evaluation of both hallucination suppression and response quality under uncertainty.

\subsubsection{Metrics}

\paragraph{Hallucination Metric}
We measure hallucination using a binary judgment indicating whether a response is supported by available information or contains unsupported or incorrect content. 
We report hallucination rate (\%) across different difficulty levels.

\paragraph{Response Quality Metrics}
We adopt an LLM-as-a-judge framework with eight evaluation dimensions, including relevance, coherence, consistency, fluency, completeness, actionability, evidence use, and conciseness. 
Most dimensions are rated on a 1--5 scale, while fluency is evaluated on a 1--3 scale. 
To ensure comparability, scores are aggregated by averaging across dimensions. 
Detailed definitions of these metrics are provided in Appendix~\ref{app:metrics}.

\subsubsection{Dataset}

We construct a controlled dataset consisting of two scenarios: (i) no reliable information (all\_wrong), where the model should abstain, and (ii) partially correct information (one\_correct), where the model should generate grounded responses. 
The dataset is derived from internal sources and is not publicly available due to privacy constraints.


\subsection{Main Results}

\subsubsection{Hallucination Reduction}

\begin{table}[h]
\centering
\caption{Hallucination rate (\%) across context sizes (lower is better).}
\label{tab:hallucination}
\begin{tabular}{cccc}
\toprule
K & GPT-4o & Base & Ours \\
\midrule
1 & 14.3\% & 24.3\% & \textbf{6.0\%} \\
3 & 16.3\% & 22.0\% & \textbf{8.3\%} \\
5 & 17.0\% & 23.3\% & \textbf{8.6\%} \\
7 & 23.0\% & 25.6\% & \textbf{8.6\%} \\
9 & 19.6\% & 26.0\% & \textbf{10.3\%} \\
\midrule
\textbf{Avg} & 18.04\% & 24.24\% & \textbf{8.36\%} \\
\bottomrule
\end{tabular}
\end{table}

As shown in Table~\ref{tab:hallucination}, our method substantially reduces hallucination compared to both the base model and GPT-4o.

On average, the hallucination rate decreases to 8.36\%, corresponding to a 65.5\% relative reduction compared to the base model and a 53.6\% reduction compared to GPT-4o.

The improvement is consistent across all context sizes, indicating robustness as contextual complexity increases.

These results suggest that explicitly modeling uncertainty and behavior enables the model to make more reliable decisions, thereby reducing hallucination in a stable and generalizable manner.

\subsubsection{Generalization Performance}

\begin{table}[h]
\centering
\caption{Hallucination rate (\%) on unseen domain dataset (lower is better).}
\label{tab:generalization}
\begin{tabular}{cccc}
\toprule
K & GPT-4o & Base & Ours \\
\midrule
1 & 2.3\% & 12.7\% & \textbf{2.3\%} \\
3 & 5.3\% & 12.3\% & \textbf{5.0\%} \\
5 & 6.7\% & 14.7\% & \textbf{5.6\%} \\
7 & 9.3\% & 18.3\% & \textbf{7.6\%} \\
9 & 8.0\% & 17.3\% & 8.6\% \\
\midrule
\textbf{Avg} & 6.32\% & 15.06\% & \textbf{5.82\%} \\
\bottomrule
\end{tabular}
\end{table}

We evaluate generalization performance under domain shift by measuring hallucination rates on unseen data.

As shown in Table~\ref{tab:generalization}, the proposed method consistently reduces hallucination compared to the base model across different values of $K$.

On average, our method achieves a hallucination rate of 5.82\%, outperforming both the base model (15.06\%) and GPT-4o (6.32\%).

Notably, performance remains competitive with GPT-4o across most settings, indicating that the learned behavior generalizes effectively beyond the training domain.

These results suggest that the proposed framework improves robustness under distribution shift while maintaining strong generalization capability.

\subsubsection{Response Quality}

\begin{table}[h]
\centering
\caption{Response quality evaluation (higher is better).}
\label{tab:response_quality}
\begin{tabular}{lcc}
\toprule
Metric & Base & Ours \\
\midrule
Relevance & 4.10 & \textbf{4.62} \\
Coherence & 3.62 & \textbf{4.42} \\
Consistency & 4.69 & \textbf{4.83} \\
Fluency & 2.86 & \textbf{2.87} \\
Completeness & 3.29 & \textbf{3.95} \\
Actionability & 2.70 & \textbf{3.32} \\
Evidence Use & 2.09 & \textbf{4.79} \\
Conciseness & 3.97 & \textbf{4.37} \\
\midrule
\textbf{Average} & 3.42 & \textbf{4.15} \\
\bottomrule
\end{tabular}
\end{table}

Response quality is evaluated across multiple dimensions following LLM-based evaluation approaches such as G-Eval~\cite{geval}.

As shown in Table~\ref{tab:response_quality}, our method consistently improves response quality across all dimensions.

The largest gain is observed in evidence use (2.09 → 4.79), indicating significantly improved grounding and transparency in generated responses.

Substantial improvements are also observed in completeness and actionability, suggesting that the model produces more informative and practically useful outputs.

Overall, the average score increases from 3.42 to 4.15, demonstrating that the proposed framework enhances response quality in a comprehensive manner.

\subsubsection{General Knowledge Preservation}

To ensure that hallucination reduction is not achieved at the expense of general capabilities,
we evaluate the model on a diverse set of benchmark datasets covering:

\begin{itemize}
    \item English knowledge (MMLU, MMLU-Pro)
    \item Chinese knowledge (CMMLU, CEVAL)
    \item Instruction following (IFEval)
    \item Mathematical reasoning (GSM8K, Math-500)
\end{itemize}

\begin{table}[h]
\centering
\caption{General knowledge evaluation (accuracy, higher is better).}
\label{tab:general_knowledge}
\begin{tabular}{l l c c}
\toprule
Domain & Dataset & Base & Ours \\
\midrule
English & MMLU     & 0.7714 & 0.7686 \\
English & MMLU-Pro & 0.6034 & 0.6000 \\
Instruction & IFEval & 0.8529 & 0.8571 \\
Chinese & CMMLU    & 0.6344 & 0.6200 \\
Chinese & CEVAL    & 0.6636 & 0.6759 \\
Math & GSM8K       & 0.9574 & 0.8842 \\
Math & Math-500    & 0.9000 & 0.8000 \\
\midrule
\textbf{Average} &  & 0.7690 & 0.7436 \\
\textbf{Average (w/o Math)} &  & 0.7051 & 0.7043 \\
\bottomrule
\end{tabular}
\\
\footnotesize{
Results are computed on a sampled subset of 3,000 instances (out of 41,512), with weighted accuracy based on dataset proportions.
}
\end{table}

The results in Table~\ref{tab:general_knowledge} show that the proposed method largely preserves general knowledge capability.

When excluding mathematical reasoning tasks, the average accuracy remains nearly unchanged (0.7051 → 0.7043), 
indicating minimal degradation in general knowledge performance.

Instruction-following ability slightly improves (IFEval: 0.8529 → 0.8571), 
suggesting that behavior-oriented training enhances controllability.

While performance decreases on mathematical reasoning tasks, 
this may be attributed to the absence of explicit optimization for multi-step reasoning in the proposed reward design.

Overall, these results demonstrate that hallucination mitigation is achieved without catastrophic forgetting \cite{EWC}, 
and the model retains its general-purpose capabilities.

\subsubsection{Summary}

Overall, the proposed method achieves a strong balance between hallucination mitigation, response quality, and general capability preservation.


\subsection{Analysis}

\subsubsection{Entropy and Overconfidence}

We analyze token-level entropy:

\[
\tilde{H} = \log_{10}(H_{\text{avg}} + \epsilon)
\]

\begin{table}[h]
\centering
\caption{Average entropy under different answerability conditions. Higher values indicate greater uncertainty.}
\label{tab:entropy}
\begin{tabular}{ccc}
\toprule
K & all\_wrong & one\_correct \\
\midrule
1 & -0.6790 & -1.5007 \\
3 & -0.8695 & -1.7862 \\
5 & -0.8462 & -1.8510 \\
7 & -0.8746 & -1.7624 \\
9 & -0.8994 & -1.7480 \\
\bottomrule
\end{tabular}
\end{table}

We analyze model uncertainty by measuring entropy under different answerability conditions.

As shown in Table~\ref{tab:entropy}, entropy is consistently higher for unanswerable cases (\textit{all\_wrong}) and lower for answerable cases (\textit{one\_correct}).

This clear separation suggests that entropy serves as an effective signal for distinguishing between answerable and unanswerable scenarios.

These results support the use of entropy as a proxy for decision boundaries, enabling the model to determine when to answer and when to abstain.

\subsubsection{Abstention Behavior}

\begin{table}[h]
\centering
\caption{Average entropy under different abstention outcomes. Higher values indicate greater uncertainty.}
\label{tab:abstention}
\begin{tabular}{ccc}
\toprule
K & Abstained & Not Abstained \\
\midrule
1 & -0.6444 & -1.1390 \\
3 & -0.7498 & -1.3853 \\
5 & -0.7753 & -1.3675 \\
7 & -0.7649 & -1.3284 \\
9 & -0.7664 & -1.5338 \\
\bottomrule
\end{tabular}
\end{table}

We analyze the relationship between entropy and abstention behavior.

As shown in Table~\ref{tab:abstention}, entropy is consistently higher for abstained cases and lower for non-abstained cases.

This suggests that correct abstention is associated with higher uncertainty, while incorrect answers tend to exhibit lower entropy, indicating overconfident predictions.

These findings support the effectiveness of the proposed entropy-based penalty in encouraging appropriate abstention behavior.
\section{Conclusion}

In this work, we revisit hallucination in large language models from a behavioral perspective. 
Rather than treating hallucination as a purely generative error, we formulate it as a decision-making problem under uncertainty, involving when to answer, how to answer, and how confidently to respond.

We propose a behavior-oriented reinforcement learning framework that integrates behavior alignment to distinguish between answering and abstaining, uncertainty modeling via entropy-based overconfidence control, response quality shaping to ensure usefulness and completeness, and length regularization to prevent degenerate generation strategies.

Through extensive experiments, we demonstrate that the proposed approach reduces hallucination by over 65\%, consistently outperforming strong baselines such as GPT-4o. 
Furthermore, the learned behavior generalizes effectively to unseen domains, significantly improves response quality across multiple dimensions, and preserves general knowledge capability without catastrophic forgetting.

These results suggest that hallucination is not merely a limitation of knowledge, but a consequence of suboptimal decision-making under uncertainty. 
By explicitly modeling behavior, confidence, and response quality, the proposed framework enables language models to produce outputs that are both reliable and useful.
\appendix

\section{Evaluation Metrics}
\label{app:metrics}

We provide detailed definitions of the response quality metrics used in the main experiments.

\subsection{Relevance}
Measures whether the response focuses on essential information without introducing irrelevant content.

\subsection{Coherence}
Evaluates logical organization and clarity of the response.

\subsection{Consistency}
Assesses whether all statements are supported by the provided information.

\subsection{Fluency}
Measures grammatical correctness and readability.

\subsection{Completeness}
Evaluates whether all key aspects of the query are addressed.

\subsection{Actionability}
Measures whether the response provides actionable guidance.

\subsection{Evidence Use}
Assesses whether the response explicitly uses supporting evidence.

\subsection{Conciseness}
Evaluates whether the response is succinct without redundancy.

\bibliographystyle{plain}
\bibliography{refs/references}


\end{document}